\documentclass[11pt,a4paper]{article}
\usepackage[margin=1in]{geometry}
\usepackage{parskip}
\usepackage{enumitem}
\usepackage{booktabs}
\usepackage{longtable}
\usepackage{xcolor}
\usepackage{titlesec}
\usepackage{array}
\usepackage[hidelinks]{hyperref}

\definecolor{navy}{HTML}{0A2A66}
\definecolor{accent}{HTML}{C0392B}
\titleformat{\section}{\color{navy}\Large\bfseries}{\thesection}{0.6em}{}
\titleformat{\subsection}{\color{navy}\large\bfseries}{\thesubsection}{0.5em}{}
\titlespacing*{\section}{0pt}{16pt}{6pt}
\newcommand{\term}[1]{\textbf{#1}}
\setlist{itemsep=2pt,topsep=3pt}
\hypersetup{linkcolor=navy}

\begin{document}

% ---------------- COVER ----------------
\begin{titlepage}
\centering
\vspace*{2cm}
{\color{navy}\Huge\bfseries Viva Examination Report\par}
\vspace{1cm}
{\LARGE Detection and Automated Mitigation of\\ Application-Layer DDoS Attacks in\\ Containerized Environments\par}
\vspace{0.6cm}
{\large\itshape Attack Focus: Slowloris --- A Closed-Loop Defense\par}
\vspace{2cm}
{\large\textbf{Course:} SE-401 --- Information Security\par}
\vspace{0.4cm}
{\large\textbf{Group Members:}\par}
{\large Summayah Mohammedy \quad Amna Khan\par}
{\large Mariha Mahmood \quad Haris Bin Nasir\par}
\vspace{2cm}
{\large A complete, self-contained study and reference document\\ prepared for the viva examination.\par}
\vfill
{\itshape All attack experiments were conducted exclusively inside an isolated\\ laboratory network against a self-hosted target. No external system was involved.\par}
\end{titlepage}

\tableofcontents
\newpage

% ---------------- EXEC SUMMARY ----------------
\section{Executive Summary (Read This First)}
In one paragraph: \emph{Slowloris} is a sneaky denial-of-service attack that
quietly opens hundreds of connections to a web server and never finishes any of
them, tying up all the server's ``workers'' until real users are locked out ---
all while using almost no bandwidth, so normal alarms never fire. This is
especially dangerous for modern apps packaged in \emph{containers} (Docker), which
have small worker pools and defenses that cannot react once an attack begins. We
built a system that \emph{watches} connections live, \emph{detects} the Slowloris
pattern using a machine-learning model (XGBoost), \emph{automatically isolates}
the attacker through the container runtime, and \emph{restores} service on its own
after the threat passes --- a complete, automatic ``closed loop.'' On a real
experiment it detected the attack with an F1 score of \textbf{0.9997}, a
false-alarm rate of \textbf{0.1\%}, restored service availability from a crash
back to \textbf{100\%} within seconds, and \textbf{never once} blocked a
legitimate slow user. The novelty is not the detector --- it is the \emph{automatic
detect$\rightarrow$isolate$\rightarrow$recover loop that runs inside the container
runtime}, together with a layer-7 feature that tells a slow attacker apart from a
slow honest user.

\textbf{How to use this document for the viva:} Part~I teaches the concepts from
zero. Parts~II--IX explain the project itself. Part~XII is a glossary of every
term. \textbf{Part~XIII is a list of likely viva questions with model answers ---
study that most.} Part~XIV is the live-demo script.

% ---------------- PART I ----------------
\section{Part I --- Foundations (Assume No Prior Knowledge)}

\subsection{What a web server does}
A \term{web server} is a program that waits for \term{clients} (browsers, apps) to
connect and ask for something over \term{HTTP}, the language of the web. For each
connection, the server assigns a \term{worker} (a thread or process) to read the
request, do the work, and send a reply. A server has a \emph{limited number of
workers} --- for example 50. If all 50 are busy, the 51st user must wait or is
refused. This limit is the resource an attacker tries to exhaust.

\subsection{How an HTTP request is structured (the key to everything)}
An HTTP request is sent as text, line by line:
\begin{itemize}
  \item A \term{request line}: e.g. \texttt{GET /page HTTP/1.1}
  \item Several \term{header} lines: e.g. \texttt{Host: example.com}
  \item A \textbf{blank line} --- this is the crucial part. The blank line
  (technically \texttt{CRLFCRLF}, two line-breaks) signals ``the headers are
  finished; now process my request.''
\end{itemize}
Until that blank line arrives, the server keeps the worker \emph{waiting}, holding
the connection open, assuming the client is just slow. \textbf{Slowloris abuses
this waiting.}

\subsection{DoS vs DDoS}
A \term{DoS} (Denial of Service) attack tries to make a service unavailable to its
real users. A \term{DDoS} (Distributed DoS) does the same from \emph{many} sources
at once. ``Availability'' is one of the three pillars of security (Confidentiality,
Integrity, Availability --- the \term{CIA triad}); DoS/DDoS attacks target
availability.

\subsection{Network-layer vs application-layer attacks}
\begin{itemize}
  \item \term{Network/volumetric attacks} flood the target with huge traffic to
  clog the pipe (e.g. a UDP flood). They are loud and easy to spot by volume.
  \item \term{Application-layer (Layer-7) attacks} are quiet. They send a small
  number of \emph{valid-looking} requests that abuse how the application behaves.
  Slowloris is Layer-7: it uses tiny traffic but exhausts a server resource.
\end{itemize}

\subsection{What Slowloris actually does (step by step)}
\begin{enumerate}
  \item The attacker opens many TCP connections to the server (e.g. 200).
  \item On each, it sends the request line and a few headers --- but \emph{not}
  the terminating blank line.
  \item Every few seconds it sends \emph{one more} harmless header line on each
  connection. This resets the server's idle timer, so the server keeps waiting.
  \item Because the request never finishes, each connection keeps its worker
  occupied \emph{forever}.
  \item With 200 held connections and only 50 workers, the pool is exhausted.
  Real users get no worker and the site appears down.
\end{enumerate}
\textbf{Analogy:} imagine 200 people at a shop's 50 service counters, each
pretending to still be deciding their order, mumbling one more word every few
seconds so the clerk can't leave. No shouting, no crowd --- but nobody else can be
served.

\subsection{What containers are, and why they matter here}
A \term{container} (via \term{Docker}) packages an application with everything it
needs so it runs the same anywhere. \term{Kubernetes} orchestrates many containers
across machines. Containers are the default way modern apps are deployed. Two facts
make Slowloris worse for them:
\begin{itemize}
  \item \textbf{Small pools:} each container is sized for modest load, so its
  worker pool is small and exhausts in \emph{seconds}.
  \item \textbf{Static defenses:} container isolation (\term{namespaces},
  \term{cgroups}, \term{network policies}) is configured in advance. It partitions
  resources but does \emph{not} watch or react to a slow attack in progress.
  \item \textbf{Multi-tenancy:} many tenants share a host, so one compromised
  container can attack its neighbours over the internal network.
\end{itemize}

% ---------------- PART II ----------------
\section{Part II --- The Problem and Motivation}

\subsection{Problem statement}
Slowloris exhausts a containerized web service's worker pool using negligible
bandwidth and traffic that looks like normal slow clients. Volume-based defenses
miss it, and static container controls cannot react while it is happening. Because
it develops over seconds-to-minutes, a human operator cannot respond in time.

\subsection{Why existing defenses fail}
\begin{itemize}
  \item \textbf{Volumetric detection} looks for high traffic --- Slowloris has
  almost none, so it is invisible to it.
  \item \textbf{Static timeouts} (e.g. Apache's \texttt{mod\_reqtimeout}) can help,
  but an attacker who trickles \emph{within} the timeout window evades them, and a
  too-aggressive timeout cuts off genuinely slow honest users.
  \item \textbf{Manual response} is too slow; by the time an alert is noticed,
  downtime has already occurred.
\end{itemize}

\subsection{Motivation}
The missing piece is not another offline detector but an \emph{automatic control
loop} that observes, decides, acts through the container runtime, and safely
reverses itself --- all in real time, with no human in the loop, and without
harming legitimate slow users.

% ---------------- PART III ----------------
\section{Part III --- The Research}

\subsection{What ``the research'' is}
This is applied security research combining \emph{machine-learning-based intrusion
detection} with \emph{automated, container-native incident response}. We study
three research questions:
\begin{itemize}
  \item \textbf{RQ1 (Accuracy):} Can we detect Slowloris accurately, and crucially,
  \emph{without} flagging slow-but-honest clients?
  \item \textbf{RQ2 (Responsiveness):} How fast does the closed loop detect,
  mitigate, and recover, and does availability actually return?
  \item \textbf{RQ3 (Cost):} What does the defense cost in overhead and false
  isolations, and how does it compare to simpler approaches?
\end{itemize}

\subsection{What others have done (related work)}
\begin{itemize}
  \item \textbf{ML detectors} for application-layer DDoS (e.g. Sharif et al.\ with
  genetic algorithms; Praseed \& Thilagam with request-pattern signatures) focus on
  \emph{detection accuracy on recorded data}.
  \item \textbf{Slow-DoS analysis} (Cambiaso et al.) defines and categorises slow
  attacks --- the basis of our threat model.
  \item \textbf{SDN mitigation} (Lukaseder et al.) acts on slow attacks at the
  network fabric using packet timing, but not inside the container runtime and
  without a layer-7 completion signal.
  \item \textbf{Datasets} (CIC-IDS2017, CIC-DDoS2019) include slow-rate subsets but
  are flow-level, so they \emph{cannot} see whether HTTP headers completed.
\end{itemize}

\subsection{The gap we fill}
Prior work overwhelmingly stops at detection on offline data. Very little
\emph{closes the loop}: detecting \emph{and} automatically mitigating a slow attack
in real time, \emph{inside a container runtime}, \emph{while protecting
slow-but-honest clients}. That combination --- layer-7 feature + automatic
mitigation + container-native --- is our contribution.

% ---------------- PART IV ----------------
\section{Part IV --- Aims and Objectives (What We Set Out to Do)}
\begin{enumerate}
  \item Build a realistic, reproducible container testbed where Slowloris
  genuinely takes a service down.
  \item Engineer features --- including a \emph{layer-7 header-completion}
  feature --- that distinguish the attack from honest slow traffic.
  \item Train and compare ML models (XGBoost vs Random Forest vs Logistic
  Regression) to detect the attack.
  \item Build an \emph{automatic} mitigator that isolates the offender through the
  Docker runtime and restores service after a cooldown.
  \item Measure not just accuracy but \emph{availability, time-to-mitigate, and
  false isolations}, and prove honest slow users are never blocked.
  \item Provide a live dashboard and a full research paper.
\end{enumerate}

% ---------------- PART V ----------------
\section{Part V --- What We Built (System Design)}

\subsection{Architecture: five cooperating containers}
All five run on one isolated Docker bridge network:
\begin{itemize}
  \item \textbf{victim} --- Apache web server with a deliberately small worker pool
  (\texttt{MaxRequestWorkers = 50}); the service we protect.
  \item \textbf{attacker} --- a lab-only Slowloris script, hard-locked to the lab
  subnet (it refuses any other target).
  \item \textbf{benign} --- generates both fast normal users \emph{and}
  slow-but-honest users (the hard negatives).
  \item \textbf{detector} --- the brain: sniffs traffic, extracts features, runs
  XGBoost, and drives the mitigator. It \emph{shares the victim's network
  namespace}, so it sees all traffic to the victim and its firewall rules take
  effect there; it also holds the Docker socket to act on containers.
  \item \textbf{dashboard} --- a live web page (\texttt{localhost:8090}) showing
  saturation, detections, and the action log in real time.
\end{itemize}

\subsection{The four-stage closed loop}
Every 2 seconds the detector runs: \textbf{Monitor $\rightarrow$ Classify
$\rightarrow$ Isolate $\rightarrow$ Restore.}
\begin{enumerate}
  \item \textbf{Monitor:} track every connection's duration, bytes/sec, packet
  timing, and whether its headers ever completed.
  \item \textbf{Classify:} score each flow with XGBoost; flag a source only if it
  holds many long-lived \emph{incomplete} flows \emph{and} the server is saturated.
  \item \textbf{Isolate:} reset the held connections, drop the source IP,
  disconnect the attacker container, and reclaim the worker pool.
  \item \textbf{Restore:} after a cooldown with no further attack, automatically
  undo everything.
\end{enumerate}

\subsection{Feature engineering (and why each feature)}
For each connection we compute the following, grouped by the invariant it captures.
\begin{longtable}{>{\raggedright}p{0.32\textwidth} p{0.60\textwidth}}
\toprule
\textbf{Feature} & \textbf{Why it matters} \\
\midrule
\endhead
\texttt{duration} & Slowloris flows are long-lived. \\
\texttt{bytes\_in}, \texttt{bytes\_per\_sec} & Attack throughput is near zero. \\
\texttt{pkt\_count} & Few packets over a long time. \\
\texttt{mean\_iat}, \texttt{std\_iat} & Inter-arrival timing; a scripted trickle is
very \emph{regular} (low variance), unlike a human. \\
\texttt{idle\_time} & Time since last byte. \\
\texttt{header\_lines} & Count of header lines seen. \\
\texttt{headers\_complete} & \textbf{The layer-7 key:} 1 if the terminating blank
line was seen. Attack = always 0; honest = eventually 1. \\
\texttt{incomplete\_ratio} & $1 -$ \texttt{headers\_complete}. \\
\texttt{small\_pkt\_ratio} & Fraction of tiny packets (the trickle). \\
\texttt{conns\_from\_src} & One attacker holds \emph{many} flows at once. \\
\bottomrule
\end{longtable}
No single feature is enough (a big upload is slow; a mobile user is bursty; a
crawler is concurrent), but their \emph{conjunction}, especially the completion
bit, is highly specific to the attack.

\subsection{The machine-learning model}
We use \term{XGBoost} (a gradient-boosted decision-tree ensemble), with
\term{Random Forest} and \term{Logistic Regression} as baselines. Trees learn
\emph{feature interactions} (e.g. ``long AND incomplete AND low-throughput''),
which is why they beat the linear baseline. We use \term{SHAP} to explain which
features drove decisions --- a check that the model relies on the intended signals,
not an artefact.

\subsection{The adaptive threshold (plain-English)}
A naive detector that fires on the model score alone would risk blocking slow
honest users. Ours fires only when \emph{both}:
\begin{itemize}
  \item the source holds at least $k=20$ flows that are model-flagged,
  \emph{incomplete}, and older than $d=8$ seconds; \emph{and}
  \item the server is genuinely under pressure (worker saturation
  $\geq 80\%$).
\end{itemize}
The 8-second duration gate is key: an honest request may be briefly mid-header but
completes in under a second, so it never accumulates. Gating on real saturation
means the system is inert while the server is healthy.

\subsection{The mitigation ladder}
Actions escalate from least to most disruptive, each logged with a timestamp:
reset the held connections $\rightarrow$ drop the source IP in the firewall
$\rightarrow$ disconnect the attacker container $\rightarrow$ reclaim the worker
pool so legitimate users regain service.

\subsection{Cooldown / restore and the grace period}
A state machine lifts the firewall drop and reconnects the container after a
cooldown with no further offence --- no permanent block to clean up manually. A
short \emph{grace period} after a reclaim prevents the surge of returning honest
users from being mistaken for a new attack.

\subsection{The live dashboard}
A small web app reads the same logs the system writes and shows, live: service
online/offline, worker saturation, attack probability, offenders isolated, an
availability chart, and the full automated action log.

% ---------------- PART VI ----------------
\section{Part VI --- How We Did It (Implementation and How to Run)}

\subsection{Technology stack}
Docker \& Docker Compose (the lab), Python (attacker, benign, detector), Scapy
(packet sniffing), XGBoost + scikit-learn (ML), SHAP (explainability), Flask (the
dashboard), and the Docker SDK + iptables (the actuator).

\subsection{Dataset collection and labelling}
We run the detector in a passive ``collect'' mode while generating fast-benign,
slow-honest, and attack traffic together. Each flow is auto-labelled by source
(attacker container $=$ attack, everything else $=$ benign), giving a clean,
automatically-labelled dataset whose negative class contains genuinely slow honest
flows. The result was about 9{,}000 labelled rows.

\subsection{How to run the live demo}
\begin{enumerate}
  \item \texttt{docker compose up -d} --- start the whole lab + dashboard.
  \item Open \texttt{http://localhost:8090} in a browser.
  \item Launch the attack:\\
  \texttt{docker compose exec -d attacker python slowloris\_lab.py --target victim
  --sockets 200 --interval 3 --duration 60}
  \item Watch the dashboard: saturation climbs, attack probability spikes, the
  offender is isolated, and availability snaps back.
  \item \texttt{docker compose down} --- stop everything.
\end{enumerate}

\subsection{How to run the full experiment (regenerate everything)}
\texttt{bash experiments/run\_all.sh} collects data, trains the models, runs the
closed-loop defense experiment, and emits every figure and table.
\texttt{python experiments/make\_results\_tex.py} then injects the real numbers
into the paper.

\subsection{How to read the results}
Outputs land in \texttt{results/}: \texttt{metrics.json} (detection scores),
\texttt{summary.json} (availability + timing), \texttt{actions.csv} (every
automated action with timestamps), \texttt{timeline.csv} (saturation over time),
\texttt{availability.csv} (user success over time), and the figures
\texttt{fig\_*.png}.

% ---------------- PART VII ----------------
\section{Part VII --- What We Achieved (Results)}

\subsection{Detection accuracy (RQ1)}
\begin{center}
\begin{tabular}{lccccc}
\toprule
Model & Precision & Recall & F1 & ROC-AUC & FPR \\
\midrule
Logistic Regression & 0.9966 & 1.00 & 0.9983 & 1.00 & 0.0062 \\
Random Forest & 0.9989 & 1.00 & 0.9994 & 1.00 & 0.0021 \\
\textbf{XGBoost} & \textbf{0.9994} & \textbf{1.00} & \textbf{0.9997} & \textbf{1.00} & \textbf{0.0010} \\
\bottomrule
\end{tabular}
\end{center}
SHAP ranked \texttt{headers\_complete} and \texttt{incomplete\_ratio} as the two
most important features --- exactly the layer-7 signals our design emphasises.

\subsection{Availability and responsiveness (RQ2)}
\begin{center}
\begin{tabular}{lcl}
\toprule
Phase & Availability & Meaning \\
\midrule
Baseline & 100\% & Healthy \\
During attack & 82\% & Brief dip; mitigation fires in $\approx$1\,s \\
After mitigation & 100\% & Pool reclaimed \\
Recovered & 100\% & Attacker auto-released \\
\bottomrule
\end{tabular}
\end{center}
Time-to-detect: within one 2-second cycle. Time-to-mitigate: about 1 second.
False isolations of benign users: \textbf{zero}. The service was down only for the
few seconds between saturation and reclaim --- not for the whole attack as in the
undefended case.

\subsection{Interpretation (what the numbers mean)}
The near-perfect detection scores reflect that \emph{once the header-completion
feature exists, the two classes are cleanly separable}. The important results are
therefore (a) the very low false-positive rate against slow-honest traffic, and
(b) the working closed loop that restores availability automatically. Accuracy
alone is not the achievement; the automatic recovery with zero collateral damage
is.

% ---------------- PART VIII ----------------
\section{Part VIII --- Novelty and Contributions}
\begin{enumerate}
  \item \textbf{A layer-7 header-completion feature} that separates a slow attacker
  from a slow honest user --- information flow-only detectors cannot see.
  \item \textbf{An adaptive threshold} that combines the ML score with live server
  saturation, so honest slow clients are safe while the server is healthy.
  \item \textbf{A closed-loop, container-native mitigator} that isolates the
  offender through the Docker runtime, reclaims the pool, and auto-restores --- the
  core contribution.
  \item \textbf{A reproducible testbed, live dashboard, and evaluation} that
  measures availability and false isolations, not just detection accuracy.
\end{enumerate}

\section{Part IX --- Significance (What It Means)}
Detection without response leaves a human in the critical path, which is too slow
for an attack that unfolds in seconds. By turning detection into an automatic,
reversible action expressed in the container runtime's own vocabulary
(``disconnect this container,'' ``restart that one''), the defense fits how modern
platforms already manage tenants and composes with existing orchestration. The
result is a service that stays available through an attack with no operator action
and no harm to legitimate users --- a practical pattern for multi-tenant container
hosting.

% ---------------- PART X ----------------
\section{Part X --- Limitations and Honest Caveats}
\begin{itemize}
  \item \textbf{Encryption:} the header-completion feature needs to read request
  bytes, so it applies to plain HTTP or behind a TLS-terminating proxy; end-to-end
  HTTPS would require terminating TLS at the monitor.
  \item \textbf{Evasion:} an attacker who randomises timing stresses the timing
  features, but not header completion, which stays discriminative.
  \item \textbf{Single host:} the testbed is one host with one attacker;
  distributed sources and cluster-wide (Kubernetes) isolation are future work.
  \item \textbf{Tuning:} the thresholds are deployment-specific and must be fitted
  to a service's pool size and traffic.
  \item \textbf{The WSL2 reclaim caveat (important --- examiners may probe this):}
  On a normal Linux server the mitigator force-closes the attacker's sockets
  instantly using \texttt{ss -K} (the kernel's \texttt{SOCK\_DESTROY} feature). Our
  test machine runs the WSL2 kernel, which does \emph{not} implement
  \texttt{SOCK\_DESTROY}, so network-disconnect alone leaves the half-open sockets
  in the kernel's backlog where fresh workers re-accept them. To reclaim reliably
  on this kernel we therefore \emph{restart the victim container}, which clears the
  backlog and brings up a fresh pool. The availability probe runs on the host, so
  the recovery it measures is real and independent of this. This is a platform
  detail, not a flaw in the approach: on production Linux the lighter socket-kill
  would be used instead.
\end{itemize}

\section{Part XI --- Future Work}
Kubernetes deployment with NetworkPolicy-based pod isolation; eBPF-based monitoring
for lower overhead and line-rate feature extraction; validation on public
benchmarks (CIC-IDS2017, CIC-DDoS2019); distinguishing genuine flash crowds from
attacks; and extension to the rest of the slow-DoS family (Slow POST, Slow Read,
HTTP/2 variants) and to TLS-terminating deployments.

% ---------------- PART XII GLOSSARY ----------------
\section{Part XII --- Glossary of Terms}
\begin{description}[leftmargin=0pt,style=nextline]
  \item[Adaptive threshold] A decision rule that combines the model score with live
  server saturation, so action is taken only when the server is genuinely
  threatened.
  \item[Application-layer (Layer-7) attack] An attack that abuses application
  behaviour (like HTTP) using valid-looking requests, rather than raw traffic
  volume.
  \item[cgroups] Linux control groups; bound a container's CPU and memory.
  \item[CIA triad] Confidentiality, Integrity, Availability --- the three core
  security goals. DoS attacks target Availability.
  \item[Closed loop] A control system that senses, decides, acts, and then observes
  the effect and adjusts --- here, detect $\rightarrow$ isolate $\rightarrow$
  restore, automatically.
  \item[Container] A lightweight, isolated package of an app and its dependencies
  (via Docker) that runs consistently anywhere.
  \item[conntrack] A Linux tool to view/remove kernel connection-tracking entries;
  used to tear down the attacker's connections.
  \item[CRLFCRLF] The two consecutive line-breaks (a blank line) that terminate an
  HTTP header block. Slowloris never sends it.
  \item[DDoS] Distributed Denial of Service --- a DoS from many sources.
  \item[DoS] Denial of Service --- making a service unavailable to its users.
  \item[Docker] The most common container platform.
  \item[eBPF] A Linux technology for safe, fast in-kernel programs; proposed for
  lower-overhead monitoring in future work.
  \item[F1 score] The harmonic mean of precision and recall; a single balanced
  measure of detection quality (1.0 is perfect).
  \item[False-positive rate (FPR)] The fraction of benign items wrongly flagged as
  attacks. Low FPR means honest users are not blocked.
  \item[Feature] A measured property of a connection (e.g.\ duration) fed to the
  model.
  \item[Flow] One connection, tracked over time (keyed by source IP and port).
  \item[Header completion] Whether an HTTP request's headers finished (the blank
  line arrived). Our key layer-7 feature.
  \item[iptables] The Linux firewall tool used to drop the attacker's IP.
  \item[Kubernetes] A system for orchestrating many containers across machines.
  \item[MaxRequestWorkers] Apache's limit on simultaneous worker processes; the
  pool Slowloris exhausts (set to 50 in our lab).
  \item[Namespaces] Linux feature that gives a container its own isolated view of
  processes, network, etc.
  \item[Network policy] A Kubernetes rule controlling which pods may talk to which;
  static and preventive.
  \item[Precision] Of the items flagged as attacks, the fraction that truly were.
  \item[Recall] Of the true attacks, the fraction that were caught.
  \item[ROC-AUC] A threshold-independent measure of how well the model separates
  the two classes (1.0 is perfect).
  \item[Saturation] Fraction of workers busy (busy / (busy + idle)); our pressure
  signal.
  \item[Scapy] A Python library for capturing and inspecting packets.
  \item[SHAP] A method that attributes a model's decision to individual features
  (explainability).
  \item[Slow-DoS family] Low-bandwidth attacks that hold a server resource:
  Slowloris (slow headers), Slow POST (slow body), Slow Read (slow response).
  \item[Slowloris] The attack studied here: many connections with never-finished
  HTTP headers, exhausting the worker pool.
  \item[SOCK\_DESTROY] A Linux kernel feature (used via \texttt{ss -K}) that
  force-closes a socket; absent on the WSL2 test kernel (see Part X).
  \item[XGBoost] A high-performance gradient-boosted decision-tree model; our main
  classifier.
\end{description}

% ---------------- PART XIII Q&A ----------------
\section{Part XIII --- Likely Viva Questions and Model Answers}
\textbf{Study this section most. Answers are written the way you would say them.}

\begin{description}[leftmargin=0pt,style=nextline]
\item[Q1. In one sentence, what is your project?]
An automatic, container-native system that detects a Slowloris slow-DoS attack in
real time using machine learning and then isolates the attacker and restores the
service on its own, without blocking legitimate slow users.

\item[Q2. What is Slowloris and why is it hard to detect?]
It opens many connections and keeps each alive with partial, never-finished HTTP
headers, exhausting the server's worker pool. It is hard to detect because it uses
almost no bandwidth and its packets look like a normal slow client, so
volume-based defenses miss it.

\item[Q3. Why is this worse in containers?]
Containers have small worker pools that exhaust in seconds, and their isolation
controls are configured in advance and cannot react to an attack in progress. In
multi-tenant hosting, one compromised container can attack its neighbours.

\item[Q4. What is the single most important idea in your work?]
The \emph{header-completion} feature combined with an \emph{automatic closed loop}.
Completion tells a slow attacker (never finishes) from a slow honest user
(eventually finishes); the loop acts on that in real time.

\item[Q5. Why machine learning and not a simple rule?]
A single rule either misses variants or flags honest slow users. The attack is
defined by a \emph{combination} of properties (long, incomplete, low-throughput,
concurrent). A tree model learns that combination, and our adaptive threshold adds
server saturation so we only act when truly threatened.

\item[Q6. Why XGBoost specifically?]
It is a strong gradient-boosted tree model that learns feature interactions and
handles tabular features well. We compared it against Random Forest and Logistic
Regression; XGBoost had the best F1 and the lowest false-positive rate.

\item[Q7. Your accuracy is 99.97\% --- isn't that suspiciously high / overfitting?]
It is high because, once the header-completion feature exists, the classes are
genuinely easy to separate --- an unfinished request is structurally different from
a finished one. We guard against a spurious result by (a) including genuinely slow
honest clients as hard negatives, (b) using a held-out test split, and (c) checking
with SHAP that the model relies on the completion features, not on something like
the attacker's IP. The real achievement is the near-zero false-positive rate and
the working recovery, not the raw accuracy.

\item[Q8. How do you avoid blocking a legitimate slow user?]
Three safeguards: we only act when the server is actually saturated ($\geq 80\%$);
we only count connections that stay \emph{incomplete} for more than 8 seconds
(honest ones finish in under a second); and we apply a grace period after mitigation
so returning users are not mistaken for a new attack. In our runs, zero honest
users were ever blocked.

\item[Q9. Walk me through what happens from attack start to recovery.]
The attack opens 200 connections and saturates the 50-worker pool within seconds.
The monitor sees many long-lived incomplete flows from one source; the classifier
flags it because saturation is high; the mitigator resets the connections, drops
the IP, disconnects the attacker container, and reclaims the pool. Availability
returns within a few seconds. After a cooldown with no further attack, the firewall
rule is lifted and the container reconnected automatically.

\item[Q10. What does ``closed loop'' mean and why does it matter?]
It means the system not only detects but also acts and then observes and reverses
itself --- sense, decide, act, recover. It matters because the attack unfolds in
seconds, far faster than a human can respond, so the response must be automatic.

\item[Q11. What is the adaptive threshold?]
Instead of acting on the model score alone, we require two things together: enough
long-lived incomplete flows from one source, \emph{and} the server being genuinely
saturated. This keeps the system inert when the server is healthy, protecting slow
honest users.

\item[Q12. How did you build your dataset and is it biased?]
We generated fast-benign, slow-honest, and attack traffic together in our lab and
auto-labelled each flow by its source container. It is narrower than a public
benchmark, which we acknowledge, but it is precisely labelled and importantly its
benign class contains real slow clients. Our methodology also admits public
datasets (CIC-IDS2017/2019) for future validation.

\item[Q13. How does mitigation actually stop the attack?]
We drop the attacker's IP in the firewall (stopping new/continued traffic),
disconnect the attacker container from the network, and reclaim the worker pool so
the seized slots are freed for real users. The attacker, now isolated, cannot
re-grab them.

\item[Q14. Why do you restart the victim container to reclaim --- isn't that heavy?]
On a normal Linux server we would force-close the attacker's sockets instantly with
\texttt{ss -K} (\texttt{SOCK\_DESTROY}). Our test kernel (WSL2) does not support
that, so the half-open sockets linger and fresh workers re-accept them. Restarting
the container reliably clears them on this platform. Because the availability probe
runs on the host, the measured recovery is real. On production Linux the lighter
socket-kill would be used and no restart would be needed.

\item[Q15. What are precision, recall, F1, and FPR?]
Precision: of what we flagged as attack, how much truly was. Recall: of all real
attacks, how many we caught. F1: their balanced average. FPR: of benign items, how
many we wrongly flagged --- the one we most want low, because it means not blocking
honest users.

\item[Q16. What is SHAP and why did you use it?]
SHAP explains which features drove each prediction. We used it to confirm the model
decides based on header completion and incompleteness --- the signals we
engineered --- rather than on an artefact of the setup.

\item[Q17. How is your work different from prior ML detectors?]
Prior work mostly stops at detection on recorded data. We close the loop with
automatic, container-native mitigation and recovery, and we use a layer-7
completion feature that flow-level detectors cannot compute.

\item[Q18. How is it different from an SDN or firewall-based defense?]
An SDN defense acts at the network fabric using timing/rate and cannot see
header completion. Ours acts inside the container runtime --- the same vocabulary
the platform uses to manage tenants --- and uses the layer-7 signal. The two are
complementary.

\item[Q19. What are the limitations?]
HTTPS hides the layer-7 feature unless TLS is terminated at the monitor;
timing-randomising attackers stress the timing features (but not completion); the
testbed is single-host; thresholds need tuning; and the reclaim mechanism is
platform-specific (the WSL2 caveat).

\item[Q20. Could the attacker evade your detector?]
They can randomise timing to blur the inter-arrival features, but they cannot
complete the request without releasing the worker --- so the completion feature
remains discriminative. If they complete requests, it is no longer Slowloris.

\item[Q21. What is the contribution if the detection is ``easy''?]
The contribution is not the detector's accuracy --- it is the automatic,
container-native loop that \emph{acts} on the detection in real time and restores
service with zero false isolations, plus the layer-7 feature that makes the
honest-vs-attacker distinction possible.

\item[Q22. How would you deploy this in production / at scale?]
Move to Kubernetes with NetworkPolicy-based pod isolation for cluster-wide scope,
use eBPF for low-overhead line-rate monitoring, terminate TLS at a proxy so the
layer-7 feature works on HTTPS, and re-fit the thresholds to the real service's
pool size and traffic profile.

\item[Q23. Is running this attack legal/ethical?]
Yes, as we did it: everything runs inside an isolated Docker network against our own
victim container, and the attack script is allow-listed to the lab subnet and
refuses any external target. We never touched a real or third-party system.

\item[Q24. What did you personally learn / what was hardest?]
The hardest part was making recovery actually work on our kernel: isolating the
attacker was not enough because half-open sockets lingered. Diagnosing that it was
the missing \texttt{SOCK\_DESTROY} kernel feature, and engineering a reliable
reclaim around it, was the key engineering lesson.

\item[Q25. If you had more time, what is the first thing you would add?]
Kubernetes deployment with pod-level isolation and validation on a public dataset,
to show the approach generalises beyond our single-host lab.
\end{description}

% ---------------- PART XIV DEMO ----------------
\section{Part XIV --- Live Demo Script}
\begin{enumerate}
  \item \textbf{Set up (before examiners watch):} \texttt{docker compose up -d};
  open \texttt{http://localhost:8090}. Point out the service is \emph{ONLINE} and
  workers are idle.
  \item \textbf{Show the baseline:} normal/slow benign traffic is flowing and the
  availability chart is at 100\%. Say: ``these include deliberately slow honest
  users --- the system must not block them.''
  \item \textbf{Launch the attack:} run the attacker command. Narrate: ``200
  connections, each sending headers that never finish.''
  \item \textbf{Watch detection:} saturation gauge climbs toward 100\%; attack
  probability spikes. Say: ``the model sees many long-lived incomplete flows from
  one source while the server is saturated.''
  \item \textbf{Watch mitigation:} the action log shows \texttt{DROP\_IP},
  \texttt{ISOLATE\_CONTAINER}, \texttt{RECLAIM\_POOL}. Availability snaps back to
  100\%.
  \item \textbf{Point out zero false positives:} the benign source was never
  touched --- only the attacker.
  \item \textbf{Watch restore:} after the cooldown the log shows
  \texttt{RESTORE\_IP} / \texttt{RESTORE\_CONTAINER}; the system is back to normal
  with no manual action.
  \item \textbf{Close:} ``this detect-isolate-reclaim-restore sequence, running
  automatically inside the container runtime, is our contribution.''
  \item \textbf{Stop:} \texttt{docker compose down}.
\end{enumerate}
\textbf{Backup:} if the live demo misbehaves, show the saved figures in
\texttt{results/} (\texttt{fig\_availability.png}, \texttt{fig\_timeline.png}) and
the paper's results tables --- they are from a real recorded run.

% ---------------- PART XV ----------------
\section{Part XV --- Team Roles (Suggested Split)}
\begin{center}
\begin{tabular}{>{\raggedright}p{0.32\textwidth} p{0.60\textwidth}}
\toprule
\textbf{Area} & \textbf{Deliverable} \\
\midrule
Lab \& victim & Docker testbed, vulnerable Apache, attack reproduction. \\
Detection & Feature extraction, model training, SHAP, baselines. \\
Mitigation & Closed loop: isolate, reclaim, cooldown/restore, dashboard. \\
Evaluation \& paper & Experiments, figures, and the write-up. \\
\bottomrule
\end{tabular}
\end{center}

\section*{Key References (full list in the paper)}
Cambiaso et al., \emph{Slow DoS Attacks: Definition and Categorisation} (2013);
Praseed \& Thilagam, \emph{HTTP Request Pattern Based Signatures} (2022);
Sharif et al., \emph{Detection of Application-layer DDoS using ML and Genetic
Algorithms} (2023); Lukaseder et al., \emph{SDN-Assisted Mitigation of Slow DDoS}
(2018); Sharafaldin et al., CIC-IDS2017 (2018); Chen \& Guestrin, \emph{XGBoost}
(2016); Lundberg \& Lee, \emph{SHAP} (2017).

\vfill
\begin{center}\itshape
End of report. All experiments were conducted inside an isolated laboratory
network against a self-hosted target.
\end{center}

\end{document}
